\documentclass[11pt, reqno]{amsart}

\usepackage[spanish]{babel}
\usepackage[LGR, T1]{fontenc}
\usepackage[utf8]{inputenc}

../LaTeX/general.tex
% \input{../graphics.tex}

\makeatletter
\def\emailaddrname{\textit{Correo electrónico}}
\def\subtitle#1{\gdef\@subtitle{#1}}
\def\@subtitle{}

% Metadata
\def\logo#1{\gdef\@logo{#1}}
\def\@logo{}
\def\institution#1{\gdef\@institution{#1}}
\def\@institution{}
\def\department#1{\gdef\@department{#1}}
\def\@department{}
\def\professor#1{\gdef\@professor{#1}}
\def\@professor{}
\def\course#1{\gdef\@course{#1}}
\def\@course{}
\def\coursecode#1{\gdef\@coursecode{#1}}
\def\@coursecode{}

\renewcommand{\maketitle}{
\begin{center}
	\small
	\renewcommand{\arraystretch}{1.2}
	\begin{tabular}{cp{.37\textwidth}p{0.44\textwidth}}
		% \hline
		\multirow{5}{*}{\includegraphics[height=2.0cm]{\@logo}}
	  & \multicolumn{2}{c}{ \makecell{{\bfseries \@institution} \\ \@department} } \\
	  % & \multicolumn{2}{c|}{{\bfseries\@institution} \\ \@department} \\
	  \cline{2-3}
	  & \textbf{Profesor:} \@professor & \textbf{Ayudante:} \authors \\
	  % \cline{2-3}
	  & \textbf{Curso:} \@course & \textbf{Sigla:} \@coursecode \\
	  % \cline{2-3}
	  & \multicolumn{2}{l}{ \textbf{Fecha:} \@date } \\
	  % \hline
	\end{tabular}
	\\[\baselineskip]
	% {}
	% \vspace{2\baselineskip}
	{\bfseries\Large\@title}
	\ifx\@subtitle\@empty\else
		\\[1ex]
		\large\mdseries\@subtitle
	\fi
\end{center}
}
\makeatother

\usepackage{multirow, makecell}

\usepackage[
	reversemp,
	letterpaper,
	% marginpar=2cm,
	% marginsep=1pt,
	margin=2.3cm
]{geometry}
\usepackage{fontawesome}
% \makeatletter
% \@reversemargintrue
% \makeatother

% Símbolos al margen, necesitan doble compilación
\newcommand{\hard}{\marginnote{\faFire}}
\newcommand{\hhard}{\marginnote{\faFire\faFire}}

% Dependencias para los teoremas
\usepackage{xifthen}
\def\@thmdep{}
\newcommand{\thmdep}[1]{
	\ifthenelse{\isempty{#1}}
	{\def\@thmdep{}}
	{\def\@thmdep{ (#1)}}
}
\newcommand{\thmstyle}{\color{thm}\sffamily\bfseries}

% ===== Estilos de Teoremas ==========
\newtheoremstyle{axiomstyle}
	{0.3cm}
	{0.3cm}
	{\normalfont}
	{0.5cm}
	{\bfseries\scshape}
	{:}
	{4pt}
	{\thmname{#1}\thmnote{ #3}\thmnumber{ (#2)}}
\newtheoremstyle{styleC}
	{0.5cm}
	{0.5cm}
	{\normalfont}
	{0.5cm}
	{\bfseries}
	{:}
	{4pt}
	{\thmname{#1\textrm{\@thmdep}}\thmnumber{ #2}\thmnote{ (#3)}}

% ====== Teoremas (sin borde) ===========
\theoremstyle{axiomstyle}
\newtheorem*{axiom}{Axioma}

% ====== Teoremas (sin borde) ==================
\theoremstyle{styleC}
\newtheorem{thm}{Teorema}[section]
\newtheorem{mydef}[thm]{Definición}
\newtheorem{prop}[thm]{Proposición}
\newtheorem{cor}[thm]{Corolario}
\newtheorem{lem}[thm]{Lema}
\newtheorem{con}[thm]{Conjetura}
\newtheorem*{sol}{Solución}
\newtheorem*{obs}{Observación}
\newtheorem*{ex}{Ejemplo}

\newtcbox{bluebox}[1][]{enhanced jigsaw, 
  sharp corners,
  frame hidden,
  nobeforeafter,
  listing only,
  #1} % comando para crear cajas de colores

\expandafter\let\expandafter\oldproof\csname\string\proof\endcsname
\let\oldendproof\endproof
\renewenvironment{proof}[1][\proofname]{%
  \oldproof[\scshape Demostración:]%
}{\oldendproof} % comando para redefinir la caja de la demostración
\newenvironment{hint}[1][\proofname]{%
  \oldproof[\scshape Pista:]%
}{\oldendproof} % comando para redefinir la caja de la demostración

% colores utilizados
\definecolor{numchap}{RGB}{249,133,29}
\definecolor{chap}{RGB}{6,129,204}
\definecolor{sec}{RGB}{204,0,0}
\definecolor{thm}{RGB}{106,176,240}
\definecolor{thmB}{RGB}{32,31,31}
\definecolor{part}{RGB}{212,66,66}

% ====== Diseño de los titulares ===============
\usepackage[explicit]{titlesec} % para personalizar el documento, la opción <<explicit>> hace que el texto de los titulares sea un objeto interactuable

\titleformat{\subsection}[runin]
	{\bfseries}
	{\textrm{\S}\thesubsection}
	{1ex}
	{#1.}

\setlist[enumerate,1]{label=\arabic*., ref=\arabic*} % Enumerate standards

\usepackage{tikz}
\usetikzlibrary{babel,cd}
% \DeclareMathOperator\Gr{Gr}

\title{Complejos de módulos}
\date{\DTMdate{2026-08-12}}

\author{José Cuevas Barrientos}
\email{josecuevasbtos@uc.cl}
\urladdr{https://josecuevas.xyz/teach/2026-2-ayud/}

\logo{../puc_negro.png}
\institution{Pontificia Universidad Católica de Chile}
\department{Facultad de Matemáticas}
\course{Topología Algebraica}
\coursecode{MAT2850}
\professor{Mauricio Bustamante}

\begin{document}

\maketitle

\section{Complejos}
Sea $R$ un anillo.
\begin{mydef}
	Un \strong{complejo de $R$-módulos} es una familia de $R$-módulos $(A_n)_{n\in\Z}$ indexada por los enteros junto a homomorfismos
	$\ud^A_n \colon A_n \to A_{n-1}$, llamados \strong{diferenciales}, tales que $\ud^A_{n-1}\circ\ud^A_n = 0$.
	De no haber ambigüedad obviaremos supra- y subíndice en el diferencial $\ud$, de modo que con abuso de notación escribimos
	\textquote{$\ud^2 = 0$}, y denotaremos solo por $A_\bullet$ al complejo.

	Un \strong{homomorfismo de complejos}, denotado $\varphi_\bullet \colon A_\bullet \to B_\bullet$, es una familia de homomorfismos
	$\varphi_n \colon A_n \to B_n$ tal que $\ud^B_n\circ\varphi_n = \varphi_{n-1}\circ\ud^A_n$ para cada $n \in \Z$.
	Es decir, tal que el siguiente diagrama conmuta:
	\[\begin{tikzcd}
		\cdots \rar & A_n \rar["\ud"] \dar["\varphi_n"] & A_{n-1} \rar["\ud"] \dar["\varphi_{n-1}"] & \cdots \\
		\cdots \rar & B_n \rar["\ud"]                   & B_{n-1} \rar["\ud"]                       & \cdots
	\end{tikzcd}\]
	Definimos el \emph{$n$-ésimo módulo de homología} de un complejo $A_\bullet$ como
	\[
		H_n(A_\bullet) := \frac{\ker( \ud \colon A_n \to A_{n-1} )}{\Img( \ud \colon A_{n+1} \to A_n )},
	\]
	un complejo $A_\bullet$ se dice \strong{exacto} o \emph{acíclico} en un índice $n \in \Z$ si $H_n(A_\bullet) = 0$;
	y se dice \emph{exacto} a secas si lo es en cada $n$.
\end{mydef}

\section{Ejercicios}
\begin{enumerate}
	\item\lookright Sea $\varphi_\bullet \colon A_\bullet \to B_\bullet$ un homomorfismo entre complejos de $R$-módulos.
		\begin{enumerate}
			\item Muestre que $\ker\varphi_\bullet$, $\Img\varphi_\bullet$ y $\coker\varphi_\bullet$ son también complejos.
			\item Muestre que induce, para cada $n \in \Z$, un homomorfismo entre módulos de homología $H_n(\varphi) \colon
				H_n(A) \to H_n(B)$.
		\end{enumerate}

	\item Sean $A_\bullet, B_\bullet$ un par de complejos de $R$-módulos.
		Dote a $(A \oplus B)_n := A_n \oplus B_n$ de un diferencial, de modo que las (familias de) inclusiones $\iota^A\colon A_\bullet \hookto (A\oplus
		B)_\bullet$, $\iota^B\colon B_\bullet \hookto (A\oplus B)_\bullet$, y las (familias de) proyecciones $\pi^A \colon (A\oplus
		B)_\bullet \to A_\bullet$ y $\pi^B \colon (A\oplus B)_\bullet \to B_\bullet$ sean homomorfismos \underline{de complejos}.
		Pruebe que existe un isomorfismo natural
		\[
			\forall n\in\Z \qquad
			H_n(A\oplus B) \cong H_n(A) \oplus H_n(B).
		\]

	\item (Prueba de lucidez)
		Sea 
		\begin{tikzcd}[cramped, sep=small]
			0 \rar & A_\bullet \rar["\varphi"] & B_\bullet \rar["\psi"] & C_\bullet \rar & 0
		\end{tikzcd}
		una sucesión exacta de complejos.
		Muestre que son equivalentes:
		\begin{enumerate}
			\item Existe un homomorfismo de complejos $\rho \colon B_\bullet \to A_\bullet$ (llamado \emph{retracción}) tal que
				$\rho\circ\varphi = \Id_{A_\bullet}$.
			\item Existe un homomorfismo de complejos $\sigma \colon C_\bullet \to B_\bullet$ (llamado \emph{sección}) tal que
				$\psi\circ\sigma = \Id_{C_\bullet}$.
			\item Hay un isomorfismo natural $B_\bullet \cong (A\oplus C)_\bullet$ como \underline{complejos}.
		\end{enumerate}
		En cuyo caso, decimos que la sucesión de complejos \emph{se escinde}.
		\begin{hint}
			Imite la prueba para módulos.
		\end{hint}

	\item Dé un ejemplo de una sucesión exacta de complejos
		\begin{tikzcd}[cramped, sep=small]
			0 \rar & A_\bullet \rar["\varphi"] & B_\bullet \rar["\psi"] & C_\bullet \rar & 0
		\end{tikzcd}
		tal que para $n \in \Z$ se escinde (es decir, $B_n \cong A_n\oplus C_n$ \textquote{de manera abstracta}), pero tal que la
		sucesión de complejos no se escinde.
		\begin{hint}
			Tal contraejemplo inclusive existe cuando el anillo es un cuerpo.
		\end{hint}
\end{enumerate}

\section{Algunos usos del lema de la serpiente}
\begin{enumerate}[resume]
	\item \textbf{Lema de los cinco:}
		Dado un diagrama conmutativo \[\begin{tikzcd}
			A^1 \rar \dar["\varphi^1"] & A^2 \rar \dar["\varphi^2"] & A^3 \rar \dar["\varphi^3"] & A^4 \rar \dar["\varphi^4"] &
			A^5 \dar["\varphi^5"] \\
			B^1 \rar                   & B^2 \rar                   & B^3 \rar                   & B^4 \rar                   &
			B^5
			\ar[from=1-1, to=2-1, draw=none, "\sim"', sloped]
			\ar[from=1-2, to=2-2, draw=none, "\sim"', sloped]
			% \ar[from=1-3, to=2-3, draw=none, "\sim"', sloped]
			\ar[from=1-4, to=2-4, draw=none, "\sim"', sloped]
			\ar[from=1-5, to=2-5, draw=none, "\sim"', sloped]
		\end{tikzcd}\]
		donde $A^\bullet$ y $B^\bullet$ son exactos, y donde $\varphi^j$ es un isomorfismo para todo $j \ne 3$.
		Pruebe que, entonces $\varphi^3$ también es un isomorfismo.
		\begin{hint}
			Puede ser útil primero probar que dado un diagrama conmutativo de módulos
			% https://q.uiver.app/#q=WzAsNCxbMCwwLCJBIl0sWzEsMCwiQSciXSxbMCwxLCJCIl0sWzEsMSwiQiciXSxbMiwzLCJcXGJldGEiLDJdLFswLDEsIlxcYWxwaGEiXSxbMCwyLCJcXHNpbSIsMl0sWzEsMywiXFxzaW0iXV0=
			\[\begin{tikzcd}
				A & {A'} \\
				B & {B'}
				\arrow["\alpha", from=1-1, to=1-2]
				\arrow[sloped, "\sim"', from=1-1, to=2-1]
				\arrow[sloped, "\sim", from=1-2, to=2-2]
				\arrow["\beta"', from=2-1, to=2-2]
			\end{tikzcd}\]
			Se tiene que $\Img\alpha \cong \Img\beta$, $\ker\alpha \cong \ker\beta$ y $\coker\alpha \cong \coker\beta$.
		\end{hint}

	\item\lookright Pruebe que una sucesión exacta corta de complejos 
		\begin{tikzcd}[cramped, sep=small]
			0 \rar & A_\bullet \rar["\varphi"] & B_\bullet \rar["\psi"] & C_\bullet \rar & 0
		\end{tikzcd}
		induce una sucesión exacta larga en homología
		\[\begin{tikzcd}
			\cdots \rar & H_{n+1}(C) \rar["\partial", red] & H_n(A) \rar["\varphi_*"] & H_n(B) \rar["\psi_*"] & H_n(C)
			\rar["\partial", red] & H_{n-1}(A) \rar & \cdots
		\end{tikzcd}\]
\end{enumerate}

\printbibliography

\end{document}
